\section*{6. Pipeline nhận dạng giọng nói tiếng Việt qua Gemini API}
\addcontentsline{toc}{section}{6. Pipeline nhận dạng giọng nói tiếng Việt qua Gemini API}

\subsection*{6.1. Lý do loại bỏ suy luận cục bộ NVIDIA/CUDA và chuyển sang kiến trúc Gemini-only}
\addcontentsline{toc}{subsection}{6.1. Lý do loại bỏ suy luận cục bộ NVIDIA/CUDA và chuyển sang kiến trúc Gemini-only}

Trong giai đoạn khảo sát ban đầu của dự án (xem hồ sơ kiểm toán lịch sử ngày 02--03/11/2026 tại \path{docs/evidence/README.md}), phương án nhận dạng giọng nói cục bộ trên GPU NVIDIA (Whisper/CUDA) đã được thử nghiệm nhưng gặp rào cản phần cứng và môi trường hệ thống:
\begin{itemize}[leftmargin=1.2cm]
  \item Lệnh kiểm tra trình điều khiển \texttt{nvidia-smi} trả về mã lỗi \texttt{9} (\textit{``NVIDIA-SMI has failed because it couldn't communicate with the NVIDIA driver''}), khiến các thư viện suy luận CUDA không thể khởi tạo trên máy trạm hiện tại.
  \item Việc cài đặt các thư viện học sâu cục bộ trực tiếp vào môi trường Python~3.10 của ROS~2 Humble tiềm ẩn xung đột ABI nghiêm trọng với các gói mở rộng C/C++ của hệ thống (đặc biệt là yêu cầu \texttt{numpy<2} của \texttt{python3-opencv} và các gói ROS~2 trên Ubuntu~22.04).
\end{itemize}

Từ quyết định kiến trúc đó, toàn bộ nhánh mã nguồn suy luận cục bộ đã được loại bỏ khỏi runtime chính và thay thế hoàn toàn bằng kiến trúc \textbf{Gemini-only} (\path{tb4_voice_control/speech.py}). Các tài liệu hoặc log kiểm thử CUDA/Whisper trước đây chỉ được giữ lại như lịch sử quyết định kỹ thuật, không dùng làm bằng chứng vận hành cho phiên bản cuối.

Để bảo vệ môi trường ROS~2 Humble, các thư viện phục vụ thu âm và gọi Gemini API được cô lập trong môi trường ảo \path{.venv-speech} (tạo bằng \texttt{python3 -m venv --system-site-packages .venv-speech}) với các ràng buộc phiên bản chặt chẽ trong \path{requirements-speech.txt}:

\begin{svcode}
google-genai>=2.33,<2.34
numpy<2
sounddevice>=0.4.6,<1.0
\end{svcode}

Phiên bản thực tế đã được xác minh khi chạy kịch bản cài đặt là \texttt{google-genai==2.33.0}, \texttt{numpy==1.21.5} và \texttt{sounddevice==0.5.6} trên Python~3.10.12.

\subsection*{6.2. Cấu hình thu âm Push-to-Talk và mã hóa WAV nội bộ trong RAM}
\addcontentsline{toc}{subsection}{6.2. Cấu hình thu âm Push-to-Talk và mã hóa WAV nội bộ trong RAM}

Thông số vận hành của bộ nhận dạng được khai báo tại \path{config/voice.yaml}:

\begin{svcode}
sample_rate: 16000
max_record_seconds: 15
provider: gemini
model: gemini-3.5-transcribe
language: vi
request_timeout_seconds: 120
\end{svcode}

Khi khởi động \texttt{voice\_control} (\path{tb4_voice_control/voice_app.py}), chương trình kiểm tra bắt buộc \texttt{provider == 'gemini'}; mọi giá trị khác đều bị từ chối ngay bằng ngoại lệ \texttt{ValueError('Speech provider must be exactly gemini')}.

Quá trình thu âm và tiền xử lý âm thanh diễn ra qua hai bước:
\begin{enumerate}[leftmargin=1.2cm]
  \item \textbf{Thu âm nhấn--giữ (\textit{Push-to-Talk}) an toàn luồng (\texttt{AudioRecorder}):}
  Khi người dùng nhấn và giữ nút \textbf{Giữ để nói} trên giao diện, tín hiệu \texttt{pressed} kích hoạt \texttt{SpeechWorker.start\_recording()} trên một \texttt{QThread} riêng. Lớp \texttt{AudioRecorder} (\path{tb4_voice_control/audio_capture.py}) chỉ nạp mô-đun \texttt{sounddevice} khi bắt đầu thu (\textit{lazy import}) và mở luồng \texttt{InputStream(samplerate=16000, channels=1, dtype='float32')}. Hàm gọi lại âm thanh (\texttt{\_audio\_callback}) sao chép từng khối mẫu đơn kênh (\texttt{mono}) vào danh sách bộ đệm dưới khóa \texttt{threading.Lock}, đồng thời giới hạn tổng số mẫu thu được không vượt quá $\text{sample\_rate} \times \text{max\_record\_seconds} = 16000 \times 15 = 240\,000$ mẫu. Khi người dùng thả nút (\texttt{released}), luồng thu được đóng an toàn và trả về một mảng liên tục \texttt{np.ndarray} kiểu \texttt{float32}.

  \item \textbf{Mã hóa PCM16 WAV trong bộ nhớ (\texttt{\_encode\_wav}):}
  Hàm \texttt{\_encode\_wav} trong \path{tb4_voice_control/speech.py} giới hạn biên độ mẫu \texttt{audio} vào đoạn $[-1, 1]$, nhân với $32767$ để lượng tử hóa sang số nguyên 16-bit có dấu thứ tự byte nhỏ trước (\texttt{<i2}), rồi đóng gói thành luồng WAV chuẩn (1 kênh, độ rộng mẫu 2 byte, tần số $16\,000\,\text{Hz}$) trực tiếp trong bộ đệm \texttt{io.BytesIO()} và mã hóa sang chuỗi ASCII Base64. Không có bất kỳ file âm thanh \texttt{.wav} tạm thời nào được ghi xuống ổ cứng, giúp bảo vệ quyền riêng tư của dữ liệu giọng nói.
\end{enumerate}

\subsection*{6.3. Giao thức gọi Generate Content API và bảo mật thông tin xác thực}
\addcontentsline{toc}{subsection}{6.3. Giao thức gọi Generate Content API và bảo mật thông tin xác thực}

Lớp \texttt{GeminiTranscriber} (\path{tb4_voice_control/speech.py}) khởi tạo đối tượng \texttt{genai.Client(api\_key=api\_key)} theo cơ chế khởi tạo trễ (\textit{lazy initialization}) ở lần nhận dạng đầu tiên. Khóa xác thực chỉ được đọc từ biến môi trường của tiến trình:
\begin{equation*}
  \texttt{api\_key = os.environ.get('GEMINI\_API\_KEY', '').strip()}
\end{equation*}
Không có khóa API nào được lưu trong file cấu hình, mã nguồn, log kiểm thử hay tài liệu báo cáo. Kịch bản \path{scripts/setup_speech_venv.sh} chỉ in trạng thái \texttt{GEMINI\_API\_KEY: present} hoặc \texttt{absent}, và \path{scripts/run_voice_control.sh} kiểm tra biến môi trường có ký tự khác khoảng trắng trước khi nạp ROS và mở giao diện.

Mỗi yêu cầu nhận dạng gửi tới phương thức \texttt{client.models.generate\_content(**request)} có cấu trúc:
\begin{itemize}[leftmargin=1.2cm]
  \item \textbf{Mô hình đích (\texttt{model}):} \texttt{gemini-3.5-transcribe}.
  \item \textbf{Nội dung (\texttt{contents}):} Gồm một lượt \texttt{role: 'user'} với hai phần (\texttt{parts}): (1) câu lệnh hướng dẫn tiếng Việt nghiêm ngặt \textit{``Chuyển lời nói tiếng Việt trong audio thành đúng một câu lệnh điều hướng. Chỉ trả về nguyên văn câu lệnh; nếu không nghe rõ, trả về chuỗi rỗng. Không giải thích.''}, và (2) khối \texttt{inline\_data} chứa dữ liệu WAV Base64 với \texttt{mime\_type: 'audio/wav'}.
  \item \textbf{Cấu hình quyền riêng tư và thời gian chờ (\texttt{config}):} Đặt \texttt{store: False} (yêu cầu không lưu trữ dữ liệu phía dịch vụ), \texttt{automatic\_continuation: False} và \texttt{http\_options: \{'timeout': request\_timeout\_seconds * 1000\}} (tương ứng $120\,000\,\text{ms}$ tức $120\,\text{s}$). Trong SDK \texttt{google-genai} 2.33.x, việc không cấu hình \texttt{retry\_options} đảm bảo tầng giao vận HTTP chỉ thực hiện đúng một lần gửi duy nhất cho mỗi thao tác thả nút, không tự động gửi lặp gây trễ ngầm.
\end{itemize}

\subsection*{6.4. Giải mã phản hồi chuyên biệt (\texttt{audio\_transcription}) và phân loại lỗi}
\addcontentsline{toc}{subsection}{6.4. Giải mã phản hồi chuyên biệt (\texttt{audio\_transcription}) và phân loại lỗi}

Một điểm kỹ thuật quan trọng đã được phát hiện và xử lý trong quá trình hoàn thiện adapter (commit \texttt{80aeed5}) là cấu trúc phản hồi của mô hình chuyên biệt \texttt{gemini-3.5-transcribe}. Khác với các mô hình sinh văn bản đa dụng trả kết quả trực tiếp qua thuộc tính \texttt{response.text}, mô hình \texttt{gemini-3.5-transcribe} trả về các đoạn văn bản nhận dạng bên trong trường \texttt{part.audio\_transcription.text} của ứng viên đầu tiên (\texttt{response.candidates[0].content.parts}).

Hàm \texttt{\_transcription\_text(response)} trong \path{tb4_voice_control/speech.py} xử lý cả hai định dạng theo thứ tự ưu tiên:
\begin{enumerate}[leftmargin=1.2cm]
  \item Duyệt tuần tự qua danh sách \texttt{response.candidates[0].content.parts}, trích xuất mọi chuỗi hợp lệ từ \texttt{part.audio\_transcription.text}, cắt khoảng trắng hai đầu (\texttt{strip()}) và ghép các đoạn không rỗng bằng đúng một dấu cách (\texttt{' '.join(...)}).
  \item Nếu danh sách \texttt{parts} không chứa trường \texttt{audio\_transcription}, hàm lùi về đọc thuộc tính \texttt{response.text.strip()} để giữ tương thích ngược.
  \item Nếu phản hồi có cấu trúc hợp lệ nhưng chuỗi văn bản rỗng (\texttt{""}, ví dụ khi micro thu khoảng lặng hoặc không nghe rõ), hàm vẫn trả về \texttt{""} hợp lệ để bộ điều khiển trung tâm báo mã \texttt{empty\_transcript}. Chỉ khi cả hai nguồn đều là \texttt{None}, hàm mới phát sinh lỗi \texttt{SpeechModelError('Speech service returned a malformed transcript.')}.
\end{enumerate}

Khi xảy ra ngoại lệ trong quá trình gọi API, hàm \texttt{\_api\_error\_message(error)} chỉ dựa vào tên lớp ngoại lệ (\texttt{\_\_mro\_\_}) và mã trạng thái HTTP (\texttt{status\_code}/\texttt{code}) để tạo thông báo hướng dẫn người dùng, tuyệt đối không in nguyên văn chuỗi lỗi từ máy chủ bên ngoài nhằm tránh rò rỉ thông tin nhạy cảm:

\begin{longtable}{|p{0.34\textwidth}|p{0.56\textwidth}|}
\hline
\textbf{Điều kiện lỗi / Mã HTTP} & \textbf{Thông báo chuẩn hóa (\texttt{SpeechModelError})} \\
\hline
\endfirsthead
\hline
\textbf{Điều kiện lỗi / Mã HTTP} & \textbf{Thông báo chuẩn hóa (\texttt{SpeechModelError})} \\
\hline
\endhead
Thiếu biến môi trường \texttt{GEMINI\_API\_KEY} & \texttt{Speech setup requires GEMINI\_API\_KEY.} \\
\hline
Lỗi khởi tạo SDK / cấu hình client & \texttt{Speech setup failed; check google-genai installation and configuration.} \\
\hline
HTTP \texttt{504}, \texttt{TimeoutError}, \texttt{TimeoutException}, \texttt{APITimeoutError} & \texttt{Speech transcription timeout; try again.} \\
\hline
HTTP \texttt{401} hoặc HTTP \texttt{403} & \texttt{Speech authentication failed; check GEMINI\_API\_KEY.} \\
\hline
HTTP \texttt{429} & \texttt{Speech rate limit reached; try again later.} \\
\hline
\texttt{ConnectionError}, \texttt{TransportError}, \texttt{APIConnectionError} & \texttt{Speech network connection failed; check connectivity.} \\
\hline
Các lỗi API/vận chuyển khác & \texttt{Speech transcription failed; try again.} \\
\hline
Phản hồi không chứa dữ liệu transcript (\texttt{None}) & \texttt{Speech service returned a malformed transcript.} \\
\hline
\end{longtable}

\reportfigure[Giao diện Voice Control UI trong bước khởi tạo, hiển thị nhãn Gemini API — gemini-3.5-transcribe, ba khu vực hợp lệ, các nút điều khiển và trạng thái đang chờ localization/action server; ảnh không để lộ API key]{03_voice_control_ui.png}{Giao diện điều khiển bằng giọng nói (\texttt{VoiceControlWindow}) trong khi chờ Nav2 và localization sẵn sàng}
